\subsection{Wartungsaufwand}
\label{subsec:wartungsaufwand}

Die Single-Repository-Strategie bündelt Core, Features, Profile, Generator,
Tests und Dokumentation. Das vermeidet parallele Baseline-Änderungen und nützt
künftigen Projekten. Bereits generierte Repositories müssen Korrekturen jedoch
eigenständig übernehmen.
Die Source of Truth reduziert somit die Zahl zu pflegender Ausgangsstände,
nicht aber den Aufwand für die Migration bestehender Produkte.

Die zentrale Pflege umfasst mehrere Sprachen, Dependency-Manager, Toolchains,
Plattformen und Releasezyklen. Diese Vielfalt folgt aus den Web-, Backend-,
Desktop- und Deploymentzielen, erhöht aber Kompetenzbedarf, Testbreite und
Umgebungsabhängigkeit.
Native Builds und externe Dienste sind dabei wartungsintensiver als ein rein
webbasierter, technologiehomogener Starter.

Fünf Profile, sechs Features und eine auswählbare Capability begrenzen derzeit
den Kombinationsraum; die Abnahme deckt fünf Basispfade und drei zulässige
PostgreSQL-Erweiterungen ab \parencite{kleiveist2026acceptance}. Weitere
Capabilities erweitern Abhängigkeiten und Testmatrix. Eine kombinatorische
Explosion ist nicht nachgewiesen, ein wachsender Pflegebedarf aber plausibel.
Katalog, Resolver, Konfigurationsvertrag, Testmatrix und Dokumentation müssen
bei solchen Erweiterungen gemeinsam konsistent bleiben.

Die zentrale Wartung ist damit aufwendiger als bei einem kleinen homogenen
Starter, kann aber wiederholte Basispflege bündeln. Ihre Wirtschaftlichkeit
hängt von Nutzungshäufigkeit, Projektähnlichkeit und Pflegezeit ab; Kostendaten
fehlen.
\beginnerinput{workspace/statement/700_bewertung/740}
